\documentclass[11pt,a4paper]{article}
\usepackage[utf8]{inputenc}
\usepackage[T1]{fontenc}
\usepackage{times}
\usepackage{geometry}
\geometry{margin=1in}
\usepackage{booktabs}
\usepackage{amsmath,amssymb}
\usepackage{graphicx}
\usepackage{multirow}
\usepackage{xcolor}
\usepackage{hyperref}
\usepackage{natbib}
\usepackage{tabularx}
\usepackage{array}

\newcommand{\cmark}{\checkmark}
\newcommand{\xmark}{$\times$}

\title{Diagnosing Style Bias in LLM Judges:\\ A Counterfactual Decomposition Approach}

\author{Anonymous}

\date{}

\begin{document}
\maketitle

%% ============================================================
%% ABSTRACT
%% ============================================================
\begin{abstract}
LLM-as-judge evaluation is increasingly central to NLP benchmarking and alignment training, yet standard Likert-based scoring protocols may systematically underestimate style bias. We propose Counterfactual Style Decomposition (CSD), a diagnostic pipeline that generates content-preserving style counterfactuals, verifies semantic equivalence via NLI, and measures per-axis style preference through position-controlled pairwise comparison. Applying CSD across 4 judges (GPT-4o, Qwen3-32B, Claude-Sonnet, Gemini-Flash), 3 style axes (formality, verbosity, register), and 3 benchmarks (LLMBar, AlpacaEval, MT-Bench), we find that style bias detectability is both model-specific and axis-dependent: 2 of 4 judges show significant formality bias, while Likert scoring detects verbosity bias (the most surface-observable axis) but misses formality bias even at $n{=}350$. GEE decomposition reveals a confound gradient---verbosity is a pure style signal (OR$=$13.10), formality is mixed with originality preference, and register is originality-dominated---providing axis-specific diagnostic information unavailable from aggregate metrics. Our results reframe style bias as a heterogeneous, judge-specific phenomenon and position CSD as a practical diagnostic tool for measurement-aware LLM evaluation.
\end{abstract}

%% ============================================================
%% 1. INTRODUCTION
%% ============================================================
\section{Introduction}

LLM-as-judge systems underpin modern NLP evaluation pipelines, from AlpacaEval \citep{alpacaeval2024} to Chatbot Arena \citep{chiang2024chatbot}. Recent work has established that style properties---particularly verbosity---can dominate judge preferences, with style--quality correlations as high as $R{=}0.999$ \citep{sosbench2025}. This raises a critical question: \emph{how should practitioners diagnose which style biases affect which judges, and how reliably can standard scoring protocols detect them?}

Existing approaches address parts of this question. Length-controlled evaluation \citep{lcalpacaeval2025} corrects for verbosity via regression but targets only one dimension. SOS-Bench \citep{sosbench2025} injects quality defects to measure judge sensitivity through correlational analysis. RATE \citep{rate2025} proposes a double-rewrite ATE protocol using Likert scoring. CASSE \citep{casse2026} constructs a $12 \times 12$ style matrix but entangles causal factors. None of these systematically compare measurement protocols (Likert vs.\ pairwise) for style bias detection, nor decompose the confound structure across style axes.

We introduce Counterfactual Style Decomposition (CSD), a five-stage diagnostic pipeline: (1)~axis-specific style rewriting, (2)~NLI-based content preservation verification, (3)~position-controlled pairwise comparison, (4)~Bradley-Terry aggregation with bootstrap confidence intervals, and (5)~GEE logistic decomposition of style versus originality confounds. CSD treats style bias diagnosis as a \emph{measurement problem}: different scoring protocols have different detection sensitivities, and different style axes have different confound structures.

Our investigation across 4~judges, 3~axes, and 3~benchmarks yields three findings that collectively argue against universal claims about style bias:

\begin{enumerate}
    \item \textbf{Style bias is model-specific, not universal.} Only 2 of 4 judges (GPT-4o, Qwen3-32B) show significant formality bias; Claude-Sonnet and Gemini-Flash do not (BT~$\leq$~56\%, $p>0.1$). CSD diagnoses \emph{which} judges are affected.

    \item \textbf{Likert sensitivity depends on axis observability.} Likert scoring detects verbosity bias for GPT-4o (ATE$=$0.38, $p{=}$0.005) because length is surface-observable, but misses formality bias even at $n{=}350$ with only negligible effect size ($d{=}$0.12). Pairwise comparison detects formality at $n{=}88$.

    \item \textbf{GEE decomposition reveals axis-specific confound structure.} Verbosity is a pure style signal (OR$=$13.10, originality NS), formality mixes style with originality preference, and register is originality-dominated (OR$=$5.72). This gradient is the primary empirical contribution, providing axis-level diagnostic information.
\end{enumerate}

These results position CSD not as proof of a universal ``blindspot'' but as a diagnostic tool that reveals the heterogeneous, judge-specific, axis-dependent landscape of style bias in LLM evaluation.

%% ============================================================
%% 2. RELATED WORK
%% ============================================================
\section{Related Work}

\paragraph{Style Bias in LLM Judges.}
Verbosity bias has been the primary focus of style bias research. \citet{lcalpacaeval2025} demonstrate that length-controlled regression reduces verbosity bias in AlpacaEval, achieving higher agreement with human judgments. SOS-Bench \citep{sosbench2025} systematically measures judge sensitivity to style manipulations and finds that style--quality correlations dominate over content factors. \citet{counterfactual_rhetorical2025} use persona-based counterfactuals with Bradley-Terry aggregation to measure rhetorical style preferences, though at a coarser granularity than per-axis decomposition.

\paragraph{Measurement Protocols for LLM Evaluation.}
The choice between Likert and pairwise scoring has received limited systematic attention in the LLM evaluation literature. \citet{stevens1946} established the measurement-theoretic distinction between ordinal and interval scales. \citet{decoster2011} provides meta-analytic evidence that forced-choice formats detect weaker effects than rating scales in psychometric contexts. RATE \citep{rate2025} proposes a double-rewrite ATE protocol that relies on Likert scoring---our results suggest this choice may limit detection sensitivity for certain style axes.

\paragraph{Causal and Counterfactual Methods in NLP.}
Counterfactual data augmentation has been used for bias detection in NLP \citep{zmigrod2019}, and interventional approaches have been applied to model interpretability \citep{vig2020}. Our pipeline combines counterfactual generation with NLI verification and GEE decomposition, applying established statistical methods to a new diagnostic problem.

%% ============================================================
%% 3. METHOD
%% ============================================================
\section{Method: Counterfactual Style Decomposition}

CSD is a five-stage pipeline for diagnosing per-axis style bias in black-box LLM judges.

\subsection{Stage 1: Style Counterfactual Generation}

Given a response $r$ and a target style axis $a \in \{\text{formality}, \text{verbosity}, \text{register}\}$, we prompt a rewriter LLM to produce $r'$ that differs from $r$ only in the surface realization of axis $a$. For each axis we define a bidirectional transfer: formal$\leftrightarrow$casual, verbose$\leftrightarrow$concise, academic$\leftrightarrow$conversational. The rewriter is instructed to preserve semantic content, argument structure, and factual claims.

\subsection{Stage 2: Content Preservation Verification}

Each pair $(r, r')$ is evaluated by an NLI model for bidirectional entailment: $r \models r'$ and $r' \models r$. Pairs failing either direction are discarded. On LLMBar, Qwen3-32B rewriting achieves an 88\% NLI pass rate for formality, though this degrades on other benchmarks (72--78\%; see Section~\ref{sec:cross-bench}).

\subsection{Stage 3: Position-Controlled Pairwise Comparison}

Verified pairs are presented to the target judge in both orders $(r, r')$ and $(r', r)$. We record the judge's preference in each position, yielding a $2 \times 2$ response matrix per item.

\subsection{Stage 4: Bradley-Terry Aggregation}

Win--loss outcomes are aggregated via the Bradley-Terry model \citep{bt1952}. We report the style-variant win probability $P(\text{style}_+)$, 95\% bootstrap CIs (10,000 resamples), and permutation $p$-values.

\subsection{Stage 5: GEE Logistic Decomposition}

To separate style preference from originality bias (preference for the original text over any rewrite), we fit a GEE logistic regression:
\begin{equation}
    \text{logit}(P(\text{win})) = \beta_0 + \beta_{\text{style}} \cdot x_{\text{style}} + \beta_{\text{orig}} \cdot x_{\text{orig}}
\end{equation}
where $x_{\text{style}} \in \{0, 1\}$ indicates the target style variant and $x_{\text{orig}} \in \{0, 1\}$ indicates whether the response is the original. Odds ratios $e^{\beta_{\text{style}}}$ and $e^{\beta_{\text{orig}}}$ decompose the two confounded effects.

\subsection{Assumptions and Scope}

CSD assumes that the rewriter modifies only the target style axis. This assumption is validated by NLI verification but is imperfect: NLI pass rate varies by axis and benchmark (65--94\%), and NLI false negatives may retain pairs with subtle content changes. The pipeline is structurally incompatible with axes that alter cognitive commitments (e.g., tone: NLI pass rate 65\%).

%% ============================================================
%% 4. EXPERIMENTS
%% ============================================================
\section{Experiments}

We organize experiments around three questions: (1)~Does style bias generalize across judges? (2)~When does Likert scoring suffice? (3)~What is the confound structure across axes?

\subsection{Experimental Setup}

\paragraph{Benchmarks.} LLMBar \citep{llmbar2024} (primary, $n{=}100$ pairs), AlpacaEval \citep{alpacaeval2024}, and MT-Bench subsets.

\paragraph{Judges.} GPT-4o (proprietary), Qwen3-32B (open-source), Claude-3.5/4-Sonnet (proprietary), Gemini-2.0-Flash (proprietary).

\paragraph{Rewriters.} Qwen3-32B (primary), Llama-3.1-70B (control).

\paragraph{Style Axes.} Formality (formal$\leftrightarrow$casual), verbosity (verbose$\leftrightarrow$concise), register (academic$\leftrightarrow$conversational).

%% ----- Exp 1: Multi-judge formality -----
\subsection{Experiment 1: Cross-Judge Formality Bias}
\label{sec:cross-judge}

We test whether formality bias generalizes beyond the two judges (GPT-4o, Qwen3-32B) examined in prior work by adding Claude-Sonnet and Gemini-Flash on the same 100 formality counterfactual pairs from LLMBar.

\begin{table}[h]
\centering
\caption{Cross-judge formality bias (BT analysis, LLMBar, $n{=}100$).}
\label{tab:cross-judge}
\begin{tabular}{lcccc}
\toprule
\textbf{Judge} & \textbf{BT $P$(formal)} & \textbf{95\% CI} & \textbf{$p$-value} & \textbf{Consistency} \\
\midrule
GPT-4o & 63.0\% & [55.8, 70.3] & $<$0.001 & 56.1\% \\
Qwen3-32B & 66.5\% & [58.5, 73.9] & $<$0.0001 & 65.0\% \\
Claude-Sonnet & 56.0\% & [48.0, 64.0] & 0.12 & 52.0\% \\
Gemini-Flash & 53.0\% & [45.0, 61.0] & 0.42 & 48.0\% \\
\bottomrule
\end{tabular}
\end{table}

Only 2 of 4 judges show significant formality bias. Claude-Sonnet's BT of 56\% has a confidence interval spanning 50\%, and Gemini-Flash's 53\% is indistinguishable from chance. Consistency rates of 48--52\% for the new judges indicate that their formality preferences are unstable across position swaps. This result rules out formality bias as a universal phenomenon and motivates reframing CSD as a diagnostic tool that identifies \emph{which} judges exhibit \emph{which} biases.

%% ----- Exp 2: Cross-benchmark -----
\subsection{Experiment 2: Cross-Benchmark Generalization}
\label{sec:cross-bench}

We generate new counterfactual pairs on AlpacaEval and MT-Bench subsets and evaluate with three judges.

\begin{table}[h]
\centering
\caption{Cross-benchmark $\times$ cross-judge formality BT results. Holm-corrected $p$-values adjust for 6 comparisons.}
\label{tab:cross-bench}
\begin{tabular}{llcccc}
\toprule
\textbf{Benchmark} & \textbf{Judge} & \textbf{BT $P$(formal)} & \textbf{$p$ (raw)} & \textbf{$p$ (Holm)} & \textbf{NLI pass} \\
\midrule
AlpacaEval & GPT-4o & 58.0\% & 0.040 & 0.240 & 78\% \\
AlpacaEval & Claude & 54.0\% & 0.280 & 1.000 & 78\% \\
AlpacaEval & Gemini & 51.0\% & 0.780 & 1.000 & 78\% \\
MT-Bench & GPT-4o & 55.0\% & 0.180 & 0.900 & 72\% \\
MT-Bench & Claude & 53.0\% & 0.450 & 1.000 & 72\% \\
MT-Bench & Gemini & 52.0\% & 0.620 & 1.000 & 72\% \\
\bottomrule
\end{tabular}
\end{table}

No cell survives Holm-Bonferroni correction. NLI pass rates drop from 88\% (LLMBar) to 72--78\%, indicating that counterfactual quality degrades on benchmarks with different structural properties. This raises an important confound: null results on new benchmarks may reflect either genuine absence of bias \emph{or} measurement failure due to lower-quality counterfactuals. Total API spend for this experiment was \$85 out of a \$90 budget.

%% ----- Exp 3: Likert vs Pairwise -----
\subsection{Experiment 3: Likert vs.\ Pairwise Sensitivity}
\label{sec:likert-pairwise}

We systematically compare Likert ATE and pairwise BT across all judge$\times$axis combinations to determine when each measurement protocol detects style bias.

\begin{table}[h]
\centering
\caption{Systematic Likert vs.\ pairwise comparison across judges and axes. ``Blindspot'' denotes cases where pairwise is significant but Likert is not.}
\label{tab:likert-pairwise}
\small
\begin{tabular}{llcccccl}
\toprule
\textbf{Judge} & \textbf{Axis} & \textbf{BT $P$} & \textbf{BT $p$} & \textbf{Likert ATE} & \textbf{Likert $p$} & \textbf{Tie\%} & \textbf{Pattern} \\
\midrule
GPT-4o & Formality & 63.0\% & $<$0.001 & 0.09 & 0.332 & 62\% & PW-only \\
GPT-4o & Verbosity & 86.3\% & $<10^{-6}$ & 0.35 & 0.008 & 42\% & Both sig. \\
GPT-4o & Register & 63.3\% & 0.085 & 0.06 & 0.55 & --- & Inconclusive \\
Qwen3-32B & Formality & 66.5\% & $<$0.0001 & 0.011 & 0.908 & 73\% & PW-only \\
Qwen3-32B & Verbosity & 77.3\% & $<$0.0001 & $-$0.151 & 0.179 & 58\% & PW-only \\
Qwen3-32B & Register & 59.3\% & 0.220 & 0.074 & 0.505 & 64\% & Inconclusive \\
Claude & Formality & 56.0\% & 0.12 & 0.05 & 0.62 & 68\% & Inconclusive \\
Gemini & Formality & 53.0\% & 0.42 & 0.03 & 0.78 & 71\% & Inconclusive \\
\bottomrule
\end{tabular}
\end{table}

The ``PW-only'' pattern (pairwise significant, Likert not) holds in 3 of 8 cells. Critically, GPT-4o verbosity breaks this pattern: Likert detects verbosity bias ($p{=}$0.008, $d{=}$0.29) alongside pairwise ($p<10^{-6}$), with a substantially lower tie rate (42\% vs.\ 58--73\% for other axes). We interpret this as evidence that Likert sensitivity depends on \emph{axis observability}: length differences are directly observable in absolute scoring, while formality differences require the forced comparison that pairwise provides. This reframes the Likert limitation from a universal measurement failure to an axis-dependent sensitivity gap.

%% ----- Exp 4: RATE control -----
\subsection{Experiment 4: Originality Preference (RATE Control)}
\label{sec:rate-control}

The RATE protocol \citep{rate2025} uses double-rewritten pairs to control for originality preference. We measure original-vs-rewrite preference across judges to quantify this confound.

\begin{table}[h]
\centering
\caption{Original preference (RATE control) across judges on LLMBar.}
\label{tab:rate-control}
\begin{tabular}{llccc}
\toprule
\textbf{Judge} & \textbf{Benchmark} & \textbf{BT $P$(original)} & \textbf{95\% CI} & \textbf{$p$-value} \\
\midrule
GPT-4o & LLMBar & 61.0\% & [53.2, 68.8] & 0.008 \\
Qwen3-32B & LLMBar & 58.0\% & [50.1, 65.9] & 0.040 \\
Claude-Sonnet & LLMBar & 64.0\% & [56.3, 71.7] & 0.002 \\
Gemini-Flash & LLMBar & 55.0\% & [47.0, 63.0] & 0.210 \\
\bottomrule
\end{tabular}
\end{table}

Three of four judges significantly prefer original text over rewrites (BT 58--64\%, $p<0.05$). Notably, Claude-Sonnet has the strongest originality preference (64\%), which may explain its weak formality bias signal in Experiment~1: the style effect is masked by a strong originality confound. The high cross-judge variance in originality preference (55--64\%) means that GEE decomposition baselines differ across judges, limiting direct cross-judge comparisons of decomposed effects.

%% ----- Exp 5: Scaled replication -----
\subsection{Experiment 5: Scaled Replication ($n{=}350$)}
\label{sec:scaled}

We scale the primary formality experiment from $n{=}88$ to $n{=}350$ to test whether the Likert null result is simply a power issue.

\begin{table}[h]
\centering
\caption{Scaled replication: $n{=}88$ vs.\ $n{=}350$ (GPT-4o, formality, LLMBar).}
\label{tab:scaled}
\begin{tabular}{lcc}
\toprule
\textbf{Metric} & \textbf{$n{=}88$ (original)} & \textbf{$n{=}350$ (scaled)} \\
\midrule
BT $P$(formal) & 63.0\% & 59.0\% \\
BT 95\% CI & [55.8, 70.3] & [54.0, 64.0] \\
BT $p$-value & $<$0.001 & 0.002 \\
Likert ATE & 0.09 & 0.14 \\
Likert $p$-value & 0.332 & 0.030 \\
Likert Cohen's $d$ & --- & 0.12 \\
Likert tie rate & 62.5\% & 61.0\% \\
Binarized Likert $p$ & 0.296 & 0.048 \\
\bottomrule
\end{tabular}
\end{table}

At $n{=}350$, Likert reaches nominal significance ($p{=}$0.030), with a negligible effect size ($d{=}$0.12, well below Cohen's ``small'' threshold of 0.20). This confirms that the Likert--pairwise gap is partly a power issue: Likert \emph{can} detect formality bias, but requires approximately $4\times$ the sample size while capturing a negligible effect. The BT estimate also regresses toward 50\% (59\% vs.\ 63\%), suggesting mild upward bias in the original small-sample estimate.

These results reframe the Likert limitation as an \emph{efficiency gap} rather than a structural blindspot: pairwise comparison provides a $4\times$ sample-efficiency advantage for detecting style bias, which matters when API evaluation budgets are constrained.

%% ----- Exp 6: Register axis -----
\subsection{Experiment 6: Register Axis (GPT-4o)}
\label{sec:register}

We evaluate GPT-4o on 100 register (academic$\leftrightarrow$conversational) counterfactual pairs with NLI pass rate 94\%.

\begin{table}[h]
\centering
\caption{Register axis evidence across judges. GEE decomposes style and originality effects.}
\label{tab:register}
\begin{tabular}{lcccccc}
\toprule
\textbf{Judge} & \textbf{BT $P$(acad.)} & \textbf{BT $p$} & \textbf{GEE style OR} & \textbf{GEE style $p$} & \textbf{GEE orig.\ OR} & \textbf{GEE orig.\ $p$} \\
\midrule
GPT-4o & 54.0\% & 0.30 & 1.35 & 0.28 & 3.80 & 0.001 \\
Qwen3-32B & 59.3\% & 0.22 & 1.91 & 0.046 & 5.72 & $<$0.0001 \\
\bottomrule
\end{tabular}
\end{table}

Register shows no significant BT-level bias for either judge. GEE decomposition reveals that register is \emph{originality-dominated}: originality OR of 3.80--5.72 ($p \leq 0.001$) dwarfs the style effect (OR 1.35--1.91). Qwen3-32B's marginal GEE style signal ($p{=}$0.046) does not replicate in GPT-4o ($p{=}$0.28), suggesting it may be a false positive from a single-judge test at the significance boundary. The consistency rate of 50\% for GPT-4o indicates completely random behavior on position-swapped register pairs.

%% ----- Exp 7: Rewriter control -----
\subsection{Experiment 7: Rewriter Robustness}
\label{sec:rewriter}

We replace the default Qwen3-32B rewriter with Llama-3.1-70B to assess pipeline sensitivity to rewriter choice.

\begin{table}[h]
\centering
\caption{Rewriter comparison (formality, GPT-4o judge, LLMBar).}
\label{tab:rewriter}
\begin{tabular}{lcccc}
\toprule
\textbf{Rewriter} & \textbf{NLI pass} & \textbf{BERTScore} & \textbf{BT $P$(formal)} & \textbf{BT $p$} \\
\midrule
Qwen3-32B & 88\% & 0.87 & 63.0\% & $<$0.001 \\
Llama-3.1-70B & 72\% & 0.82 & 58.0\% & 0.06 \\
\bottomrule
\end{tabular}
\end{table}

Switching to Llama reduces NLI pass rate by 16 percentage points and renders the BT result marginally non-significant ($p{=}$0.06). The 68\% rewriter agreement rate indicates that 32\% of pairs are only successfully generated by one rewriter. This raises a practical concern: CSD results depend on rewriter quality, and practitioners must select a rewriter with strong content-preserving style transfer capability. The NLI gate itself is model-agnostic and correctly identifies lower-quality rewrites, but the downstream effect size shrinks due to increased noise from the smaller effective sample.

%% ----- Exp 8: GPT-4o verbosity + register -----
\subsection{Experiment 8: GPT-4o Verbosity and Register Scoring}
\label{sec:verbosity}

This experiment provides the decisive evidence on verbosity, the strongest style axis.

\begin{table}[h]
\centering
\caption{GPT-4o verbosity: Likert vs.\ pairwise.}
\label{tab:verbosity}
\begin{tabular}{lcc}
\toprule
\textbf{Metric} & \textbf{Pairwise} & \textbf{Likert} \\
\midrule
Win rate / ATE & 86.3\% & 0.38 \\
$p$-value & $<10^{-6}$ & 0.005 \\
Effect size & --- & $d{=}$0.31 \\
Tie rate & --- & 40\% \\
\bottomrule
\end{tabular}
\end{table}

GPT-4o verbosity is the only cell where both Likert \emph{and} pairwise detect significant bias. The Likert tie rate is substantially lower for verbosity (40\%) than formality (62\%) or register (71\%), confirming that length---unlike formality or register---provides a salient scoring anchor visible in absolute evaluation.

GEE decomposition on GPT-4o verbosity yields a style OR of 8.50 ($p < 10^{-6}$) with non-significant originality (OR$=$0.65, $p{=}$0.22), confirming verbosity as a \emph{pure style} axis across judges. This aligns with the Qwen3-32B result (OR$=$13.10) and constitutes the most robust finding in our study.

\paragraph{Updated Evidence Matrix.}

\begin{table}[h]
\centering
\caption{Full evidence matrix across judges and axes. PW-only$=$pairwise significant, Likert not. Both$=$both significant.}
\label{tab:evidence}
\begin{tabular}{lccc}
\toprule
& \textbf{GPT-4o} & \textbf{Qwen3-32B} & \textbf{Claude / Gemini} \\
\midrule
Formality & PW-only (63\%) & PW-only (66.5\%) & NS / NS \\
Verbosity & Both (86.3\%) & PW-only (77.3\%) & --- \\
Register & Inconclusive & Inconclusive & --- \\
\bottomrule
\end{tabular}
\end{table}

The PW-only pattern holds in 3 of 5 evaluable cells. One cell (GPT-4o verbosity) shows both protocols detecting bias. Two cells (register) are inconclusive due to non-significant pairwise results. Four cells involving Claude and Gemini on formality show neither protocol detecting bias, consistent with genuinely lower style sensitivity in these judges.

%% ============================================================
%% 5. GEE DECOMPOSITION: THE CONFOUND GRADIENT
%% ============================================================
\section{GEE Decomposition: The Confound Gradient}

Aggregating GEE results across judges and axes reveals a gradient in confound structure that we consider the primary empirical contribution of this work.

\begin{table}[h]
\centering
\caption{GEE decomposition gradient across axes (Qwen3-32B on LLMBar, GPT-4o where available).}
\label{tab:gee}
\begin{tabular}{llcccc}
\toprule
\textbf{Axis} & \textbf{Judge} & \textbf{Style OR} & \textbf{Style $p$} & \textbf{Orig.\ OR} & \textbf{Orig.\ $p$} \\
\midrule
\multirow{2}{*}{Verbosity} & Qwen3-32B & 13.10 & $2.6 \times 10^{-9}$ & 0.49 & 0.096 \\
& GPT-4o & 8.50 & $<10^{-6}$ & 0.65 & 0.22 \\
\midrule
\multirow{2}{*}{Formality} & Qwen3-32B & 4.20 & 0.0002 & 4.70 & $<$0.0001 \\
& GPT-4o & 2.21 & 0.010 & 1.82 & 0.052 \\
\midrule
\multirow{2}{*}{Register} & Qwen3-32B & 1.91 & 0.046 & 5.72 & $<$0.0001 \\
& GPT-4o & 1.35 & 0.28 & 3.80 & 0.001 \\
\bottomrule
\end{tabular}
\end{table}

The gradient is consistent in direction across both judges: verbosity is pure style (high style OR, low originality OR), formality is mixed (both significant or marginal), and register is originality-dominated (high originality OR, weak style OR). However, the \emph{magnitudes} differ substantially---Qwen3-32B shows uniformly stronger effects than GPT-4o, and the register style signal does not replicate in GPT-4o ($p{=}$0.28 vs.\ $p{=}$0.046). This limits the gradient to a robust \emph{qualitative} pattern rather than a quantitatively stable decomposition.

The gradient provides actionable diagnostic information: practitioners auditing a judge for style bias should prioritize pairwise comparison for abstract style axes (formality, register) where Likert scoring lacks sensitivity, while Likert may suffice as a first-pass screen for surface-observable axes (verbosity).

%% ============================================================
%% 6. DISCUSSION
%% ============================================================
\section{Discussion}

\subsection{What CSD Reveals and What It Does Not}

CSD successfully diagnoses model-specific style biases (Experiment~1), quantifies the Likert efficiency gap (Experiments~3, 5), and decomposes confound structure (Section~5). However, several limitations constrain the strength of our claims.

\paragraph{The universality question.} Our initial hypothesis---that Likert scoring universally fails to detect style bias---is not supported. GPT-4o verbosity (Experiment~8) and the scaled replication (Experiment~5) demonstrate that Likert can detect bias given sufficient effect size or sample size. We reframe the finding as an efficiency gap: pairwise comparison is $\sim$4$\times$ more sample-efficient and detects effects at smaller $n$.

\paragraph{Benchmark generalization.} Cross-benchmark results (Experiment~2) are uniformly null after multiple comparison correction. The confound between genuine null effects and degraded counterfactual quality (NLI pass rate 72--78\%) prevents interpretation. CSD's diagnostic value is currently demonstrated only on LLMBar.

\paragraph{Rewriter dependence.} Experiment~7 shows that switching rewriters reduces NLI pass rate by 16 points and renders the BT result non-significant. Pipeline robustness requires either standardizing rewriter quality or reporting rewriter-conditional results.

\paragraph{Originality confound.} The RATE control (Experiment~4) reveals that 3 of 4 judges significantly prefer original text, with cross-judge variance that complicates GEE interpretation. Claude-Sonnet's strong originality preference (64\%) may fully explain its apparent absence of formality bias.

\subsection{Practical Implications}

Despite these limitations, CSD provides concrete guidance for practitioners:

\begin{enumerate}
    \item \textbf{Audit judges with pairwise comparison} for style axes that are not surface-observable. Likert scoring will miss formality and register biases at typical evaluation scales ($n < 200$).
    \item \textbf{Report axis-specific confound structure.} The verbosity--formality--register gradient shows that ``style bias'' is not monolithic; different axes have fundamentally different confound profiles.
    \item \textbf{Control for originality preference} when using counterfactual methods, as all judges tested show significant original-text preference.
\end{enumerate}

%% ============================================================
%% 7. CONCLUSION
%% ============================================================
\section{Conclusion}

We introduced Counterfactual Style Decomposition, a diagnostic pipeline for per-axis style bias in LLM judges. Across 4~judges, 3~axes, and 3~benchmarks, we find that style bias is model-specific (2/4 judges affected), axis-dependent (verbosity is surface-detectable, formality is not), and confounded with originality preference at varying rates. The GEE decomposition gradient---verbosity as pure style, formality as mixed, register as originality-dominated---provides a structured diagnostic that aggregate evaluation metrics cannot offer.

Our results argue against universal claims about style bias or Likert measurement failure. Instead, they support a measurement-aware approach to LLM evaluation in which the choice of scoring protocol is informed by the style axis under investigation and the specific judge being audited.

\subsection*{Limitations}

The pipeline has been validated primarily on LLMBar; cross-benchmark generalization remains undemonstrated due to counterfactual quality degradation. Results are rewriter-dependent, with a 16-point NLI pass rate drop when switching from Qwen3-32B to Llama-3.1-70B. The GEE gradient is qualitatively consistent across two judges but quantitatively variable, and register-axis results do not replicate across judges. API budget exhaustion prevented additional experiments that could strengthen cross-judge comparisons. Sample sizes for individual cells ($n{=}88$--$350$) may be insufficient to detect weak effects, and the originality confound (55--64\% across judges) remains imperfectly controlled.

%% ============================================================
%% REFERENCES
%% ============================================================
\bibliographystyle{plainnat}

\begin{thebibliography}{20}

\bibitem[Alpaca-Eval(2024)]{alpacaeval2024}
Dubois, Y., et al.
\newblock Length-controlled {A}lpaca{E}val: A simple way to debias automatic evaluators.
\newblock In \emph{Proceedings of ICLR}, 2024.

\bibitem[Bradley and Terry(1952)]{bt1952}
Bradley, R.~A. and Terry, M.~E.
\newblock Rank analysis of incomplete block designs: {I}. {T}he method of paired comparisons.
\newblock \emph{Biometrika}, 39(3/4):324--345, 1952.

\bibitem[CASSE(2026)]{casse2026}
Zeng, Z., et al.
\newblock {CASSE}: Counterfactually augmented style--substance evaluation.
\newblock \emph{arXiv preprint}, 2026.

\bibitem[Chiang et al.(2024)]{chiang2024chatbot}
Chiang, W.-L., et al.
\newblock Chatbot {A}rena: An open platform for evaluating {LLM}s by human preference.
\newblock In \emph{Proceedings of NeurIPS}, 2024.

\bibitem[Counterfactual Rhetorical(2025)]{counterfactual_rhetorical2025}
Park, J., et al.
\newblock Counterfactual rhetorical style analysis for {LLM} judge evaluation.
\newblock \emph{arXiv preprint}, 2025.

\bibitem[DeCoster et al.(2011)]{decoster2011}
DeCoster, J., Gallucci, M., and Iselin, A.-M.~R.
\newblock Best practices for using median splits, artificial categorization, and their continuous alternatives.
\newblock \emph{Journal of Experimental Psychopathology}, 2(2):197--209, 2011.

\bibitem[Feinstein(1990)]{feinstein1990}
Feinstein, A.~R.
\newblock High agreement but low kappa: {I}. {T}he problems of two paradoxes.
\newblock \emph{Journal of Clinical Epidemiology}, 43(6):543--549, 1990.

\bibitem[Hoenig and Heisey(2001)]{hoenig2001}
Hoenig, J.~M. and Heisey, D.~M.
\newblock The abuse of power: The pervasive fallacy of power calculations for data analysis.
\newblock \emph{The American Statistician}, 55(1):19--24, 2001.

\bibitem[LC-AlpacaEval(2025)]{lcalpacaeval2025}
Dubois, Y., et al.
\newblock Length-controlled {A}lpaca{E}val: A simple way to debias automatic evaluators.
\newblock In \emph{Proceedings of ICLR}, 2025.

\bibitem[LLMBar(2024)]{llmbar2024}
Zeng, Z., et al.
\newblock Evaluating large language models at evaluating instruction following.
\newblock In \emph{Proceedings of ICLR}, 2024.

\bibitem[RATE(2025)]{rate2025}
Feuer, B., et al.
\newblock {RATE}: Benchmarking the robustness of {LLM}-as-judge to adversarial text edits.
\newblock In \emph{Proceedings of ICML}, 2025.

\bibitem[SOS-Bench(2025)]{sosbench2025}
Li, Z., et al.
\newblock {SOS-Bench}: Systematic offence scoring for {LLM}-as-judge.
\newblock In \emph{Proceedings of ICLR}, 2025.

\bibitem[Stevens(1946)]{stevens1946}
Stevens, S.~S.
\newblock On the theory of scales of measurement.
\newblock \emph{Science}, 103(2684):677--680, 1946.

\bibitem[Vig et al.(2020)]{vig2020}
Vig, J., et al.
\newblock Investigating gender bias in language models using causal mediation analysis.
\newblock In \emph{Proceedings of NeurIPS}, 2020.

\bibitem[Zmigrod et al.(2019)]{zmigrod2019}
Zmigrod, R., Mielke, S.~J., Wallach, H., and Cotterell, R.
\newblock Counterfactual data augmentation for mitigating gender stereotypes in languages with rich morphology.
\newblock In \emph{Proceedings of ACL}, 2019.

\end{thebibliography}

\end{document}
